% !TEX root = tarasov-qasp-zh.tex
\cleardoublepage
\thispagestyle{empty}
\begin{figure*}[!ht]
\centering
\includegraphics[width=0.75\textwidth]{figs/sec/tarasov-qasp-sec-02.pdf}
\end{figure*}
\begin{fullwidth}
{\textsf{\Large \hlred{力的概念是基本的物理概念之一。你能足够熟练地运用它吗？你对动力学定律理解得好吗？}}}
\end{fullwidth}
\cleardoublepage


\chapter[你能画出作用在物体上的力吗？]{你能画出作用在物体上的力吗？\label{ch-02}}

\begin{dialogue}

\std 力学习题似乎是所有习题中最难的。解这类题该从哪里入手呢？

\tchr 常常可以从分析作用在物体上的力入手。比如，可以取下面几种情形（\figr{fig-07}）：
\begin{enumerate}[label=(\alph*), leftmargin=1cm]
\item 物体被斜向上抛出（与水平方向成一定角度），
\item 物体沿斜面下滑，
\item 物体在竖直平面内绕绳的一端做圆周运动，
\item 物体是一个摆。
\end{enumerate}
\begin{figure}[!h]%
\centering
\includegraphics[width=0.38\textwidth]{figs/ch-02/fig-007a.pdf}
\caption[物体的各种不同情形。]{物体的各种不同情形。每种情形都要找出作用在该物体上的力。\label{fig-07}}
\end{figure}
请对每种情形画出表示作用在物体上的力的箭头，并说明各箭头代表什么。

\std 这是我画的图（\figr{fig-08}）。第一种情形，$P$ 是物体的重力，$F$ 是抛力；第二种情形，$P$ 是重力，$F$ 是使物体沿斜面下滑的力，$F_{f\!\!r}$ 是摩擦力；第三种情形，$P$ 是重力，$F_{c}$ 是向心力（centripetal force），$T$ 是绳的张力；第四种情形，$P$ 是重力，$F$ 是回复力（restoring force），$T$ 是绳的张力。

\begin{figure}[!h]%
\centering
\includegraphics[width=0.4\textwidth]{figs/ch-02/fig-008a.pdf}
\caption[学生对图 7 各情形中物体所受作用力的画法。]{学生对\figr{fig-07} 各情形中物体所受作用力的画法。\label{fig-08}}
% 注：短题内禁用 \ref/\figr（进入 .lof 引发 label 错误），写纯文字图号
\end{figure}

\tchr 四种情形你都画错了。这里是我画的正确图示（\figr{fig-09}）。必须清楚理解的一点是：力是物体之间相互作用的结果。因此，要表示作用在物体上的力，必须先弄清有哪些物体与该物体相互作用。比如第一种情形，只有地球在吸引物体，与它相互作用（\figr{fig-09}~\drkgry{(a)}）。所以物体只受一个力——重力 $P$。如果要把空气阻力或者风的作用考虑进来，那就得引入额外的力。\hlblue{你图中所画的“抛力”其实并不存在，因为没有产生这种力的相互作用。}
\pnote{“抛力”（连同“冲力”“惯力”一类说法）是中国学生同样爱画的力。检验标准只有一条：每一个力都必须说得出施力物体——人教版必修第一册第三章开篇即强调力是物体对物体的作用。说不出施力物体的“力”，一律不存在。}

\begin{figure}[!h]%
\centering
\includegraphics[width=0.42\textwidth]{figs/ch-02/fig-009a.pdf}
\caption[每种情形下物体所受作用力的正确画法。]{每种情形下物体所受作用力的正确画法。请与\figr{fig-08} 对照。\label{fig-09}}
\end{figure}

\std 可是，要把物体抛出去，总得对它施加某种力吧。

\tchr 对，这没错。抛物体时你的确对它施加了力。但上面讨论的是物体被抛出之后的运动，也就是说，是使物体获得一定飞行初速度的那个力停止作用之后的运动。\hlblue{力是不可能“储存”起来的；相互作用一结束，力也就不存在了。}

\std 可是，如果物体只受重力作用，它为什么不沿竖直方向下落，而是沿曲线运动呢？

\tchr 在这种情形里，物体运动的方向与所受力的方向不一致，这让你感到奇怪。其实这完全符合牛顿第二定律。你的问题说明你对动力学中的牛顿定律想得还不够透彻。这个问题我打算放到后面再谈（见\chap{ch-04}）。现在我想继续分析这四种情形。第二种情形（\figr{fig-09}~\drkgry{(b)}）是物体沿斜面下滑。有哪些物体与它相互作用？

\std 显然有两个：地球和斜面。

\tchr 正是。这样就能找出作用在物体上的力了。地球给出重力 $P$；斜面则给出滑动摩擦力（sliding friction）$F_{f\!\!r}$ 和通常称为支承力（bearing reaction）的力 $N$\sidenote{也称法向力（normal force）。}。注意，你在图里把力 $N$ 完全漏掉了。
\pnote{人教版必修第一册第三章称这个力为支持力，记作 $F_N$（即原书及图中的 $N$）；重力人教版记作 $G$（原书用 $P$）。支持力属于弹力，方向垂直于接触面——“法向力”一名即由此而来。}

\std 慢着！这么说，斜面对物体作用的力是两个，而不是一个？

\tchr 当然只有一个力。不过，把它分解为两个分力来处理更方便：一个沿斜面方向（滑动摩擦力），一个垂直于斜面（支承力）。这两个力同源，即它们是同一个力的分量，这一点体现在 $F_{f\!\!r}$ 与 $N$ 之间存在一个普遍关系：
\begin{equation}%
F_{f\!\!r}=kN
\label{eq-05}
\end{equation}
其中 $k$ 是一个常量，叫做滑动摩擦系数（coefficient of sliding friction）。这个关系我们后面还要详细讨论（\chap{ch-03}）。
\pnote{式~\eqref{eq-05} 即滑动摩擦定律，人教版必修第一册第三章记作 $F=\mu F_N$，比例系数 $\mu$ 称动摩擦因数，即原书的 $k$。符号不同，内容相同。}

\std 我在图里画了一个使物体沿斜面下滑的“下滑力”。看来这种力是没有的。可我清楚地记得，过去常听人说起“下滑力”这个词。这该怎么看？

\tchr 是的，这个说法确实存在。但要记住：\hlblue{你所说的“下滑力”，不过是重力的一个分量——把重力沿斜面方向和垂直斜面方向分解得到的两个分力之一。}在列举物体所受的力时，既然已经列了重力，就没有理由再添加作为其分量的“下滑力”。第三种情形（\figr{fig-09}~\drkgry{(c)}）是物体在竖直平面内做圆周运动。有哪些物体对它作用？
\pnote{“下滑力”在中文教辅里同样流行。现行人教版不使用这个词：把重力沿斜面与垂直斜面两个方向分解，得到的是重力的两个分力，不是新增的力。判断方法仍是一条——“下滑力”说不出施力物体。}

\std 两个：地球和绳。

\tchr 好，所以作用在物体上的是两个力：重力和绳的张力。

\std 那向心力呢？

\tchr 别急！圆周运动的问题里错误太多了，我打算在后面专门详细讨论（\chap{ch-08}）。这里只想指出：\hlblue{向心力不是额外作用在物体上的又一种力，它是合力。}在我们这个情形里（物体位于轨迹最低点时），向心力等于绳的张力和重力之差。
\pnote{向心力（以及上面第四种情形的回复力）是按效果命名的力，不指某种特定的相互作用。人教版必修第二册第六章讨论圆周运动时同样强调：受力图上只画真实的力（重力、张力、摩擦力……），向心力是它们合成的效果，绝不能作为独立的箭头与它们并列——这是受力分析的常见陷阱。}

\std 要是我理解得对，第四种情形（\figr{fig-09}~\drkgry{(d)}）里的回复力也是绳的张力和重力的合力吧？

\tchr 完全正确。这里与第三种情形一样，与物体相互作用的是绳和地球。因此，作用在物体上的是两个力：绳的张力和重力。我要再强调一遍：\hlblue{力只能因物体的相互作用而产生，不会从任何“附加”的考虑中冒出来。找出与研究对象作用的物体，也就找出了作用在它上面的力。}
\pnote{原书此处从简：严格地说，摆的回复力是重力沿圆弧切线方向的分力；张力恒沿半径、与速度垂直，对回复不起作用，张力与重力的合力在一般位置并不指向平衡位置（其径向分量是向心力）。人教版选择性必修第一册第二章处理单摆即取切向分力为回复力；而教师此处的要点——回复力不是新的相互作用力——完全正确。}

\std 比您在\figr{fig-07} 中所举更复杂的情形肯定还有。可以考虑一下它们吗？

\tchr 更复杂的相互作用的例子很多。比如，一个恒定的水平力 $F$ 作用在物体上，使物体沿斜面向上运动。这种情形下物体所受的力如\figr{fig-10} 所示。

\begin{figure}
\centering
\includegraphics[width=0.4\textwidth]{figs/ch-02/fig-010a.pdf}
\caption[斜面上的物体及其所受的力。]{斜面上的物体及其所受的力。\label{fig-10}}
\end{figure}

再举一例：带电的摆悬挂在平行板电容器中做振动。这时多出一个力 $F_{e}$，是电容器电场作用在摆上电荷的力（\figr{fig-11}）。解题时可能遇到的各种情形，显然无法一一列举。

\begin{figure}[!h]%
\centering
\includegraphics[width=0.4\textwidth]{figs/ch-02/fig-011a.pdf}
\caption[平行板电容器中带电摆所受的力。]{平行板电容器中带电摆所受的力。\label{fig-11}}
\end{figure}
\std 如果题里有好几个物体，该怎么办？比如\figr{fig-12} 所示的情形。

\begin{figure}[!h]%
\centering
\includegraphics[width=0.4\textwidth]{figs/ch-02/fig-012a.pdf}
\caption[静置于斜面上的三个物体各受哪些力？]{静置于斜面上的三个物体各受哪些力？\label{fig-12}}
\end{figure}
\pnote{底本此图图题原文为 ``three bodies on balanced on an incline''，多衍一个 ``on''，显系录入之误，中文版径改。本章另有若干标点级排印小误（如 ``the, plane''、``t.he''），径改不注。}

\tchr 每一次都应当明确：你要考虑的是哪些物体、或是哪个物体组合的运动。就拿你提出的例子中物体 1 的运动来说。与它相互作用的是地球、斜面和绳 $AB$。

\std 物体 2 不与物体 1 相互作用吗？

\tchr 只通过绳 $AB$ 作用。作用在物体 1 上的力是：重力 $P'$、滑动摩擦力 $F_{f\!\!r}$、支承力 $N'$ 和绳 $AB$ 的张力 $T'$（\figr{fig-13}~\drkgry{(a)}）。

\std 可为什么您图中的摩擦力方向朝左？让它朝相反方向似乎也同样合理。

\begin{figure}[!h]%
\centering
\includegraphics[width=0.45\textwidth]{figs/ch-02/fig-013a.pdf}
\caption[如图 12 所示静置于斜面上的三个物体所受的力。]{如图~\ref{fig-12} 所示静置于斜面上的三个物体所受的力。\label{fig-13}}
\end{figure}

\tchr \hlblue{要确定摩擦力的方向，必须知道物体运动的方向。}如果题中没有指明，就应当假定其中一个方向。在本题中，我假定物体 1（连同整个物体系统）向右运动，滑轮顺时针转动。当然，我事先无法知道对不对；只有代入具体数值后，运动方向才成为确定的。如果我的假定错了，算出的加速度就会是负值。那时就得改设物体向左运动（滑轮逆时针转动），并把滑动摩擦力的方向相应地反过来。然后重新列出计算加速度的方程，再代入数值检验其符号。

\std 为什么还要再检验一次加速度的符号？若假定向右运动时它是负的，那么第二次假定下它显然就是正的了。

\tchr 不，第二种假定下它同样可能是负的。

\std 这我就不明白了。物体不向右运动，显然就得向左运动，这不明摆着吗？

\tchr 你忘了物体还可能静止。这个问题我们后面还会回来，详细分析考虑摩擦力时产生的种种复杂情况（见\chap{ch-07}）。眼下先假定滑轮顺时针转动，来看物体 2 的运动。

\std 与物体 2 相互作用的是地球、斜面、绳 $AB$ 和绳 $CD$。作用在物体 2 上的力如\figr{fig-13}~\drkgry{(b)} 所示。

\tchr 很好。现在转到物体 3。

\std 物体 3 只与地球和绳 $CD$ 相互作用。\figr{fig-13}~\drkgry{(c)} 画出了作用在物体 3 上的力。

\tchr 好。确定了每个物体所受的力之后，就可以对每个物体列出运动方程，然后求解所得的方程组。

\std 您刚才说过，不必每个物体分开处理，也可以把物体组当作一个整体来考虑。

\tchr 是啊。物体 1、2 和 3 既可以像刚才那样分开研究，也可以作为整体来研究。这时绳中的张力就不必考虑了，因为它们在这种情形下变成了内力（internal force），即所研究对象各部分之间的相互作用力。三个物体组成的系统，整体只与地球和斜面相互作用。
\pnote{这段对话展示的正是中国教辅所说的“整体法”（把物体 1、2、3 作为整体考虑，绳中张力成为内力，不出现在方程中）与“隔离法”（逐个物体列方程，可求出张力）。两者通常配合使用：先用整体法求加速度，再隔离求内力。}

\std 我想弄清一点。我在画\figr{fig-13}~\drkgry{(b)} 和~\drkgry{(c)} 中的力时，曾假定绳 $CD$ 在滑轮两侧的张力相同。这样假定对吗？

\tchr 严格地说，不对。如果滑轮在顺时针转动，绳 $CD$ 连接物体 3 的那一侧张力应当大于连接物体 2 的那一侧。正是这个张力差使滑轮加速转动。本题的例子中已假定滑轮质量可以忽略。换句话说，滑轮没有需要加速的质量，它只是被看作改变连接物体 2 和 3 的绳的方向的一种手段。因此可以认为绳 $CD$ 在滑轮两侧张力相同。\hlblue{按照惯例，只要题目没有特别说明，滑轮质量都忽略不计。}问题都弄清楚了吗？
\pnote{这就是中学熟知的“轻绳、轻滑轮”模型：不计滑轮质量（及轴摩擦）时，跨过滑轮的绳两侧张力相等。大学普通物理计入滑轮的转动惯量后，两侧张力不再相等，其差值提供滑轮加速转动所需的力矩——即上文“张力差使滑轮加速转动”的定量形式。}

\std 我还有一个关于力的作用点的问题。您在图中把所有力都画在物体的同一点上。这正确吗？比如，能把摩擦力画在物体的重心上吗？

\tchr 应当记住：我们研究的是质点（material point，或称 particle）——而不是有一定大小和形状的物体——的运动学和动力学，也就是说，把物体看作一个有质量的点。图上画出物体而不画出点，只是为了清晰。因此，\hlblue{所有力都可以画成作用在物体的同一点上。}
\pnote{质点是人教版必修第一册第一章引入的第一个理想模型：当物体的大小和形状对所研究问题的影响可以忽略时，就可把物体当作质点。不能忽略的情形，正是下面讨论的转动与翻倒。}

\std 我们学过，任何简化都会使问题的某些方面丢失。把物体看作质点，丢掉的究竟是什么？

\begin{figure}[!h]%
\centering
\includegraphics[width=0.4\textwidth]{figs/ch-02/fig-014a.pdf}
\caption[作用于 A、B 两点的力 F₁ 与 F₂（a）；将 F₂ 平移至点 A 后须附加一个力偶（b）。]{作用于 $A$、$B$ 两点的力 $F_{1}$ 与 $F_{2}$（a）；将 $F_{2}$ 平移至点 $A$ 后须附加一个力偶（$F_{2}$ 与 $F_{2}''$）（b）。\label{fig-14}}
\end{figure}
\tchr 在这种简化下，我们没有考虑转动矩（rotational moment）——在真实条件下，它可能引起物体的转动和翻倒。质点只有平动。来看一个例子。设物体上两个不同的点各作用着一个力：点 $A$ 处为 $F_{1}$，点 $B$ 处为 $F_{2}$，如\figr{fig-14}~\drkgry{(a)} 所示。现在在点 $A$ 再加上与力 $F_{2}$ 等大且平行的力 $F_{2}'$，以及与力 $F_{2}$ 等大但方向相反的力 $F_{2}''$（见\figr{fig-14}~\drkgry{(b)}）。由于力 $F_{2}'$ 与 $F_{2}''$ 相互平衡，它们的加入丝毫不改变问题的物理面貌。不过，\figr{fig-14}~\drkgry{(b)} 可以这样来解释：作用在点 $A$ 的力 $F_{1}$ 与 $F_{2}'$ 引起物体的平动；此外，物体上还作用着一个由力 $F_{2}$ 与 $F_{2}''$ 组成的力偶（force couple），引起转动。换句话说，可以把力 $F_{2}$ 平移到物体上的点 $A$，只要同时补上相应的转动矩。当我们把物体看作质点时，显然就没有转动矩了。
\pnote{底本\figr{fig-14} 的图题误沿了\figr{fig-13} 的文字（“三个物体在斜面上……”），与本图内容（力的平移与力偶）不符，显系录入之误；中文版按图面内容拟写图题，待考原书。同段底本另有排印脱漏（``cause motion of translation of the body also applied to the body is a force couple''），句意不连贯，按上下文补足为“……引起物体的平动；此外还作用着一个力偶，引起转动”，中文版照此译出。}
\pnote{力偶与力偶矩在中学不作要求。力偶是两个等大、反向、不共线的力，其特点是合力为零而合力矩不为零，只改变物体的转动状态；“把一个力平移到另一点、同时附加一个转动矩”即理论力学中的力线平移定理，是大学普通物理处理刚体问题的基本工具。质点模型丢掉的正是这类转动效应。}

\std 您说质点不能转动，只有平动。可我们已经处理过转动了——圆周运动就是。

\tchr 不要把完全不同的两回事混为一谈。点的平动可以沿各种轨迹进行，比如沿圆周。我排除质点转动的可能，指的是绕自身（即绕通过该点的任一轴）的转动。

\end{dialogue}
